% ============================================================
%  Signal Flow — Prototipo 1: Modelador de Recintos 3D
%  Compilar con: pdflatex signal_flow.tex
%  (requiere pdflatex o lualatex; tikz, pgf, amsmath, geometry)
% ============================================================
\documentclass[11pt,a4paper]{article}

\usepackage[utf8]{inputenc}
\usepackage[T1]{fontenc}
\usepackage[spanish]{babel}
\usepackage{geometry}
\geometry{margin=2.5cm}

\usepackage{amsmath, amssymb}
\usepackage{tikz}
\usetikzlibrary{arrows.meta, positioning, fit, shapes.geometric,
                shapes.multipart, calc, backgrounds, decorations.pathreplacing}
\usepackage{xcolor}
\usepackage{booktabs}
\usepackage{array}
\usepackage{enumitem}
\usepackage{hyperref}
\usepackage{caption}
\usepackage{subcaption}
\usepackage{float}

% ---- Colores del tema ----
\definecolor{clrGUI}{RGB}{137,180,250}       % azul claro — módulos GUI
\definecolor{clrGeo}{RGB}{166,227,161}       % verde — geometría
\definecolor{clrFEM}{RGB}{245,194,231}       % rosa — FEM
\definecolor{clrViz}{RGB}{250,179,135}       % naranja — visualización
\definecolor{clrMat}{RGB}{203,166,247}       % violeta — materiales
\definecolor{clrSig}{RGB}{148,226,213}       % teal — señales Qt
\definecolor{clrDark}{RGB}{30,30,46}         % fondo oscuro

% ---- Estilo de cajas ----
\tikzset{
  box/.style    = {draw, rectangle, rounded corners=4pt,
                   minimum width=3.2cm, minimum height=0.9cm,
                   text centered, align=center,
                   font=\footnotesize\bfseries,
                   inner sep=4pt},
  guibox/.style  = {box, fill=clrGUI!40},
  geobox/.style  = {box, fill=clrGeo!50},
  fembox/.style  = {box, fill=clrFEM!50},
  vizbox/.style  = {box, fill=clrViz!50},
  matbox/.style  = {box, fill=clrMat!40},
  sigbox/.style  = {box, fill=clrSig!40, dashed},
  smallbox/.style= {box, minimum width=2.4cm, minimum height=0.7cm,
                    font=\scriptsize\bfseries},
  arr/.style    = {-{Stealth[length=6pt]}, thick},
  darr/.style   = {-{Stealth[length=6pt]}, thick, dashed, gray},
  label/.style  = {font=\scriptsize, midway, fill=white, inner sep=1pt,
                   align=center},
}

% ============================================================
\begin{document}
% ============================================================

\title{\textbf{Flujo de señales}\\
       Prototipo 1 — Modelador de Recintos 3D con Simulación Acústica FEM}
\author{Ingeniería acústica / Desarrollo de software}
\date{}
\maketitle
\tableofcontents
\newpage

% ------------------------------------------------------------
\section{Visión general de la arquitectura}
% ------------------------------------------------------------

La aplicación está organizada en tres capas: \textbf{interfaz (GUI)},
\textbf{cómputo} y \textbf{renderizado 3D}.  La comunicación entre capas
ocurre mediante el sistema de \emph{signals/slots} de Qt5; ningún módulo
llama directamente a otro sin pasar por una señal o un método público
explícito.

\begin{figure}[H]
\centering
\begin{tikzpicture}[node distance=1.4cm and 2.2cm]

  % --- Capa GUI ---
  \node[guibox, minimum width=4cm] (main)    {MainWindow\\{\scriptsize main.py}};
  \node[guibox, below=0.7cm of main, xshift=-2.8cm] (ctrl)
        {ControlPanel\\{\scriptsize controls.py}};
  \node[guibox, below=0.7cm of main, xshift= 2.8cm] (apanel)
        {AcousticPanel\\{\scriptsize acoustic\_panel.py}};

  % --- Capa cómputo ---
  \node[geobox, below=1.5cm of ctrl] (geo)
        {geometry.py\\{\scriptsize make\_room / make\_arch\_ribs}};
  \node[fembox, below=1.5cm of apanel] (fem)
        {acoustic\_fem.py\\{\scriptsize build\_KM / solve\_modes / FRF}};
  \node[matbox, right=1.3cm of fem] (mlib)
        {material\_library.py\\{\scriptsize Sabine / RT60 / $\xi_n$}};
  \node[fembox, below=0.9cm of fem, xshift=-2cm] (amesh)
        {acoustic\_mesh.py\\{\scriptsize build\_volume\_mesh}};
  \node[fembox, below=0.9cm of fem, xshift= 2cm] (aana)
        {acoustic\_analysis.py\\{\scriptsize pipelines / slices / campos}};

  % --- Capa renderizado ---
  \node[vizbox, below=1.5cm of geo, xshift=2cm] (viewer)
        {IsoViewer\\{\scriptsize viewer.py}};
  \node[vizbox, right=1.2cm of viewer] (aviz)
        {acoustic\_viewer.py\\{\scriptsize markers / slices / nubes}};

  % --- Flechas ---
  \draw[arr] (main) -- node[label, left]{parámetros} (ctrl);
  \draw[arr] (main) -- node[label, right]{get\_surface} (apanel);
  \draw[arr] (ctrl) -- node[label, left]{verts/tris/edges} (geo);
  \draw[arr] (geo)  -- node[label, below]{update\_geometry} (viewer);
  \draw[arr] (apanel) -- node[label, right]{run\_fem\_modal} (fem);
  \draw[arr] (fem)  -- node[label, left]{eigsh} (amesh);
  \draw[arr] (fem)  -- node[label, right]{ModalSolution} (aana);
  \draw[arr] (mlib) -- node[label, right]{$\xi_n$} (fem);
  \draw[arr] (aana) -- node[label, below]{slices / nubes} (aviz);
  \draw[arr] (aviz) -- (viewer);

  % --- Señales Qt (punteadas) ---
  \draw[darr] (ctrl)   to[bend left=15]
      node[label, left]{\footnotesize parametersChanged} (main);
  \draw[darr] (viewer) to[bend right=20]
      node[label, below]{\footnotesize sourceAddRequested\\sourceMoveRequested}
      (main);
  \draw[darr] (main)   to[bend right=10]
      node[label, right]{\footnotesize add\_source\_at} (apanel);

\end{tikzpicture}
\caption{Arquitectura de tres capas. Líneas sólidas = llamadas directas;
         líneas punteadas = señales Qt.}
\end{figure}

% ------------------------------------------------------------
\section{Flujo 1 — Cambio de geometría}
% ------------------------------------------------------------

Cada vez que el usuario mueve un slider o edita el polígono de planta,
\texttt{ControlPanel} emite \texttt{parametersChanged}.  El flujo completo
hasta el renderizado se describe a continuación.

\begin{figure}[H]
\centering
\begin{tikzpicture}[node distance=0.9cm and 0cm]

  \node[guibox]  (s1) {Usuario mueve slider\\o edita polígono};
  \node[sigbox,  below=of s1] (s2)
        {\texttt{parametersChanged(params:dict)}\\señal Qt};
  \node[guibox,  below=of s2] (s3)
        {\texttt{MainWindow.\_on\_params()}};
  \node[geobox,  below=of s3] (s4)
        {\texttt{make\_room(**params)}\\$\rightarrow$ verts, tris, edges};
  \node[geobox,  below=of s4] (s5)
        {\texttt{make\_arch\_ribs(**params)}\\(solo si arch\_height $>$ 0)};
  \node[vizbox,  below=of s5] (s6)
        {\texttt{viewer.update\_geometry(v,t,e)}\\refresh mesh + aristas};
  \node[fembox,  below=of s6] (s7)
        {\texttt{acoustic.on\_geometry\_changed()}\\invalida modal\_result,
         $\xi_n$, campo 3D};
  \node[vizbox,  below=of s7] (s8)
        {\texttt{acoustic.\_refresh\_sources\_list()}\\re-agrega esferas
         \emph{después} del mesh};
  \node[vizbox,  below=of s8] (s9)
        {\texttt{acoustic.\_refresh\_receiver\_marker()}\\re-agrega cruz};

  \foreach \i/\j in {s1/s2, s2/s3, s3/s4, s4/s5, s5/s6, s6/s7, s7/s8, s8/s9}
    \draw[arr] (\i) -- (\j);

  % Nota lateral
  \node[right=2cm of s7, font=\scriptsize, align=left,
        draw=gray, rounded corners, fill=yellow!10, inner sep=4pt,
        minimum width=4.2cm]
    (nota1) {Invalidación necesaria:\\la malla FEM antigua\\
             no corresponde al\\nuevo recinto.};
  \draw[darr] (s7) -- (nota1);

\end{tikzpicture}
\caption{Flujo de cambio de geometría.}
\end{figure}

\subsection*{Detalles del mallado volumétrico}

\texttt{build\_volume\_mesh} construye una malla tetraédrica interior:

\begin{enumerate}[itemsep=2pt]
  \item Calcular AABB (bounding box) del recinto.
  \item Generar rejilla estructurada de hexaedros con paso $h = 1/\texttt{n\_per\_meter}$.
  \item Subdividir cada hexaedro en 6 tetraedros (split de Freudenthal).
  \item Filtrar tetraedros cuyo centroide cae \emph{fuera} del recinto
        (test point-in-polyhedron por rayo en $+z$, paridad de intersecciones).
  \item Re-indexar nodos usados.
\end{enumerate}

% ------------------------------------------------------------
\section{Flujo 2 — Cálculo modal FEM}
% ------------------------------------------------------------

\begin{figure}[H]
\centering
\begin{tikzpicture}[node distance=0.85cm and 0cm]

  \node[guibox] (f1) {Usuario presiona\\``Calcular modos (FEM)''};
  \node[fembox, below=of f1] (f2)
        {\texttt{\_solve\_fem()}};
  \node[geobox, below=of f2] (f3)
        {\texttt{get\_surface()} $\rightarrow$ (verts, tris)};
  \node[fembox, below=of f3] (f4)
        {\texttt{build\_volume\_mesh(verts, tris, n\_per\_meter)}\\
         $\rightarrow$ nodes~$(N_n\times3)$, tets~$(N_e\times4)$};
  \node[fembox, below=of f4] (f5)
        {\texttt{build\_KM(nodes, tets)}\\
         $\rightarrow$ $\mathbf{K},\mathbf{M}$ sparse CSR};
  \node[fembox, below=of f5] (f6)
        {\texttt{solve\_modes(K, M, n\_modes)}\\
         \texttt{eigsh} shift-invert (Lanczos)\\
         $\rightarrow$ $f_n$~[Hz],\; $\boldsymbol{\Phi}$~$(N_n\times N_m)$};
  \node[matbox, below=of f6] (f7)
        {\texttt{\_compute\_xi\_from\_materials()}\\
         Sabine $\rightarrow$ RT60$(f)$ $\rightarrow$ $\xi_n$};
  \node[vizbox, below=of f7] (f8)
        {\texttt{\_refresh\_modes\_combo()}\\llena ComboBox con $f_n$};
  \node[vizbox, below=of f8] (f9)
        {\texttt{\_update\_slice()}\\slice 2D del modo~0};

  \foreach \i/\j in {f1/f2, f2/f3, f3/f4, f4/f5, f5/f6, f6/f7, f7/f8, f8/f9}
    \draw[arr] (\i) -- (\j);

\end{tikzpicture}
\caption{Flujo del cálculo modal FEM.}
\end{figure}

\subsection*{Formulación variacional}

El problema de autovalores acústico con condición de Neumann homogénea
(paredes rígidas):

\[
  \mathbf{K}\,\boldsymbol{\phi}_n \;=\; \lambda_n\,\mathbf{M}\,\boldsymbol{\phi}_n,
  \qquad
  f_n = \frac{c}{2\pi}\sqrt{\lambda_n}
\]

donde
\[
  K_{ij} = \int_\Omega \nabla N_i \cdot \nabla N_j \,\mathrm{d}V,
  \qquad
  M_{ij} = \int_\Omega N_i \, N_j \,\mathrm{d}V
\]

Los modos se M-ortonormalizan: $\boldsymbol{\phi}_n^T \mathbf{M} \boldsymbol{\phi}_m = \delta_{nm}$.

\subsection*{Amortiguamiento modal desde materiales}

\[
  A(f) = \sum_i \alpha_i(f)\,S_i
  \qquad\text{(absorción de Sabine)}
\]
\[
  \mathrm{RT}_{60}(f) = \frac{0{,}161\,V}{A(f)}
  \qquad\text{(Sabine)}
\]
\[
  \xi_n = \frac{1{,}1}{f_n \cdot \mathrm{RT}_{60}(f_n)}
  \qquad\text{(amortiguamiento modal del modo } n\text{)}
\]

Los coeficientes $\alpha_i(f)$ se obtienen de archivos JSON en la carpeta
\texttt{materials/} e interpolan en escala logarítmica entre bandas de octava.

% ------------------------------------------------------------
\section{Flujo 3 — Evaluación de la FRF}
% ------------------------------------------------------------

\begin{figure}[H]
\centering
\begin{tikzpicture}[node distance=0.85cm and 0cm]

  \node[guibox] (r1) {Usuario presiona ``Calcular FRF''};
  \node[fembox, below=of r1] (r2)
        {\texttt{\_compute\_frf()}};
  \node[fembox, below=of r2] (r3)
        {\texttt{run\_fem\_frf(modal, sources, receiver,}\\
         \texttt{f\_min, f\_max, n\_freqs, damping=$\xi_n$)}};
  \node[fembox, below=of r3] (r4)
        {\texttt{acoustic\_fem.frequency\_response(...)}\\
         evalúa $H(f_k)$ para cada $f_k$};
  \node[vizbox, below=of r4] (r5)
        {\texttt{FRFDialog(result, modal\_freqs)}\\
         gráfico matplotlib embebido en Qt};
  \node[vizbox, below=of r5] (r6)
        {dB SPL $= 20\log_{10}(|H|/20\,\mu\text{Pa})$\\
         + líneas verticales en $f_n$};

  \foreach \i/\j in {r1/r2, r2/r3, r3/r4, r4/r5, r5/r6}
    \draw[arr] (\i) -- (\j);

  % Nota: export
  \node[right=1.8cm of r6, font=\scriptsize, align=left,
        draw=gray, rounded corners, fill=yellow!10, inner sep=4pt,
        minimum width=3.8cm]
    (nota2) {Export PNG/SVG/PDF\\desde botón o barra\\de herramientas.};
  \draw[darr] (r6) -- (nota2);

\end{tikzpicture}
\caption{Flujo de cálculo y visualización de la FRF.}
\end{figure}

\subsection*{Superposición modal}

\[
  H(f) = i\omega\rho_0 \sum_{n=1}^{N_m}
         \frac{\varphi_n(\mathbf{x}_r)
               \displaystyle\sum_s Q_s\,\varphi_n(\mathbf{x}_s)}
              {\omega_n^2 - \omega^2 + 2i\,\xi_n\,\omega_n\,\omega}
\]

con $\omega = 2\pi f$, $\rho_0 = 1{,}21\ \text{kg/m}^3$.
El denominador se evalúa en cada frecuencia $f_k$ del eje de la FRF;
el numerador (independiente de $\omega$) se pre-calcula.

% ------------------------------------------------------------
\section{Flujo 4 — Visualización del campo acústico 3D}
% ------------------------------------------------------------

El botón ``Actualizar campo 3D'' (o la tecla Enter) despacha
\texttt{\_update\_field\_3d()}, cuyo comportamiento depende del
\texttt{combo\_field}:

\begin{figure}[H]
\centering
\begin{tikzpicture}[node distance=0.85cm and 0cm]

  \node[guibox] (c1) {Trigger:\\botón / Enter / debounce 350~ms};
  \node[fembox, below=of c1] (c2)
        {\texttt{\_update\_field\_3d()}};

  % Bifurcación
  \node[below=0.85cm of c2] (fork) {};
  \node[geobox, below=0.6cm of c2, xshift=-4cm] (c3a)
        {\texttt{mode\_shape\_field\_3d}\\
         evalúa $\varphi_n(\mathbf{x})$ en rejilla 3D\\
         $\rightarrow$ valores con signo $[-1,1]$};
  \node[fembox, below=0.6cm of c2, xshift= 4cm] (c3b)
        {\texttt{pressure\_field\_3d}\\
         evalúa $|p(\mathbf{x},f_n)|$ en rejilla 3D\\
         (superposición modal, con fuentes)};

  \node[vizbox, below=0.85cm of c3a] (c4a)
        {\texttt{pressure\_3d.update\_signed()}\\
         Azul $\leftrightarrow$ Blanco $\leftrightarrow$ Rojo\\
         (valores negativos / cero / positivos)};
  \node[vizbox, below=0.85cm of c3b] (c4b)
        {\texttt{pressure\_3d.update()}\\
         Azul $\to$ Verde $\to$ Rojo\\
         (mínimo $\to$ máximo presión)};

  \node[vizbox, below=0.85cm of c4a, xshift=4cm] (c5)
        {\texttt{GLScatterPlotItem} (pxMode=True)\\
         montado sobre \texttt{IsoViewer}};

  \draw[arr] (c1) -- (c2);
  \draw[arr] (c2) -| node[label, above left]
      {\scriptsize combo\_field = 0\\(Forma modal)}
      (c3a);
  \draw[arr] (c2) -| node[label, above right]
      {\scriptsize combo\_field = 1\\(Presión $|p|$)}
      (c3b);
  \draw[arr] (c3a) -- (c4a);
  \draw[arr] (c3b) -- (c4b);
  \draw[arr] (c4a) |- (c5);
  \draw[arr] (c4b) |- (c5);

\end{tikzpicture}
\caption{Bifurcación del flujo de visualización según \texttt{combo\_field}.}
\end{figure}

\subsection*{Rejilla de evaluación}

Los campos se evalúan sobre una rejilla uniforme 3D dentro del AABB de la
malla FEM.  Solo se conservan los puntos donde el interpolador baricéntrico
devuelve un valor finito (puntos dentro de algún tetraedro):

\[
  \{(x_i, y_i, z_i)\}_{i=1}^N
  = \bigl\{ \mathbf{x} \in \mathbb{R}^3 \mid
    \mathrm{FEM{-}locator}(\mathbf{x}) \neq \mathrm{NaN} \bigr\}
\]

\subsection*{Por qué forma modal y presión se parecen \emph{en resonancia}}

En la frecuencia $f = f_n$ (resonancia del modo $n$):

\[
  p(\mathbf{x}) \approx
  \frac{\rho_0}{2\xi_n\,\omega_n} \cdot
  \varphi_n(\mathbf{x}_s) \cdot \varphi_n(\mathbf{x})
\]

La distribución espacial de $|p(\mathbf{x})|$ es proporcional a $|\varphi_n(\mathbf{x})|$.
La diferencia observable es el \emph{signo}: la forma modal muestra lóbulos
$+$ y $-$ (azul/rojo), mientras que $|p|$ es siempre no-negativo (solo rojo).

Al mover la fuente a un \textbf{nodo} del modo $n$ ($\varphi_n(\mathbf{x}_s)=0$),
la contribución del modo $n$ se anula y el patrón cambia significativamente
(dominan otros modos fuera de resonancia).

% ------------------------------------------------------------
\section{Flujo 5 — Posicionamiento y arrastre de fuentes}
% ------------------------------------------------------------

\begin{figure}[H]
\centering
\begin{tikzpicture}[node distance=0.8cm and 0cm]

  % --- Colocar fuente ---
  \node[guibox] (p1) {Ctrl + Click derecho\\en \texttt{IsoViewer}};
  \node[sigbox,  below=of p1] (p2)
        {\texttt{mousePressEvent}\\raycast al plano $z=0$};
  \node[sigbox,  below=of p2] (p3)
        {\texttt{sourceAddRequested(x,y,1.0)}\\señal Qt};
  \node[guibox,  below=of p3] (p4)
        {\texttt{MainWindow.\_on\_source\_add\_from\_viewer}\\
         $\to$ \texttt{acoustic.add\_source\_at(x,y,1)}};
  \node[fembox,  below=of p4] (p5)
        {\texttt{SourceArray.add(OmniSource)}\\
         actualiza lista interna};
  \node[vizbox,  below=of p5] (p6)
        {\texttt{src\_markers.update(sources)}\\
         \texttt{GLScatterPlotItem} re-creado};

  \foreach \i/\j in {p1/p2, p2/p3, p3/p4, p4/p5, p5/p6}
    \draw[arr] (\i) -- (\j);

  % --- Arrastrar fuente ---
  \node[right=3.5cm of p1, guibox] (q1) {Shift + Click izquierdo\\sobre esfera de fuente};
  \node[sigbox,  below=of q1] (q2)
        {\texttt{\_pick\_source(px,py)}\\retorna idx $\geq 0$ (fuente) ó $-2$ (receptor)};
  \node[sigbox,  below=of q2] (q3)
        {\texttt{\_dragging\_source\_idx = idx}\\guarda $z$ de arrastre};
  \node[sigbox,  below=of q3] (q4)
        {Mouse move $\to$ \texttt{\_pick\_horizontal\_plane}\\
         $\to$ \texttt{sourceMoveRequested(idx,x,y,z)}};
  \node[guibox,  below=of q4] (q5)
        {\texttt{\_on\_source\_moved\_from\_viewer}\\
         \texttt{srcs.sources[idx] = OmniSource(...)}};
  \node[vizbox,  below=of q5] (q6)
        {\texttt{src\_markers.set\_positions(srcs)}\\actualiza en-lugar (sin re-crear GL)};
  \node[fembox,  below=of q6] (q7)
        {\texttt{acoustic.schedule\_field\_update()}\\
         reinicia timer debounce 350~ms};
  \node[sigbox,  below=of q7] (q8)
        {Timer dispara $\to$ \texttt{\_deferred\_field\_update}\\
         (solo si \texttt{combo\_field}=Presión~$|p|$)};

  \foreach \i/\j in {q1/q2, q2/q3, q3/q4, q4/q5, q5/q6, q6/q7, q7/q8}
    \draw[arr] (\i) -- (\j);

  % Nota receptor
  \node[below=0.4cm of q8, font=\scriptsize, align=left,
        draw=gray, rounded corners, fill=yellow!10, inner sep=4pt,
        minimum width=4cm]
    (nota3) {Para el receptor (idx=$-2$):\\señal \texttt{receiverMoveRequested}\\
             $\to$ actualiza spinboxes\\
             $\to$ \texttt{rcv\_marker.update()}};
  \draw[darr] (q8.south) -- (nota3);

\end{tikzpicture}
\caption{Flujo de posicionamiento (izquierda) y arrastre (derecha) de fuentes acústicas.}
\end{figure}

\subsection*{Raycast para posicionamiento}

El click en el visor 3D lanza un rayo desde el ojo de la cámara en la
dirección del pixel seleccionado.  La intersección con el plano $z = z_0$
se calcula como:

\[
  t = \frac{z_0 - o_z}{d_z}, \qquad
  \mathbf{x} = \mathbf{o} + t\,\mathbf{d}
\]

donde $\mathbf{o}$ es el origen del rayo y $\mathbf{d}$ su dirección,
obtenidos invirtiendo la matriz de proyección-vista construida manualmente
con \texttt{QMatrix4x4}.

\subsection*{Detección de fuentes en pantalla (\textit{picking})}

Para Shift+Click, la función \texttt{\_pick\_source(px,py)} proyecta cada
posición de fuente $(x_s,y_s,z_s)$ al espacio de pantalla y busca la más
cercana dentro de un radio de 28~px:

\[
  i^* = \arg\min_{i} \bigl[(px - u_i)^2 + (py - v_i)^2\bigr],
  \quad \text{si} < 28^2 \text{ px}^2
\]

donde $(u_i, v_i)$ es la proyección de la fuente $i$.  El receptor usa el
índice especial $-2$.

% ------------------------------------------------------------
\section{Flujo 6 — Slice 2D horizontal}
% ------------------------------------------------------------

\begin{figure}[H]
\centering
\begin{tikzpicture}[node distance=0.8cm and 0cm]

  \node[guibox] (sl1)
    {Cambio en:\\combo\_mode / combo\_field / sb\_slice\_z / sb\_slice\_res};
  \node[fembox, below=of sl1] (sl2)
    {\texttt{\_update\_slice()}};

  \node[fembox, below=0.8cm of sl2, xshift=-3.5cm] (sl3a)
    {\texttt{slice\_mode\_shape}\\$\varphi_n(x,y,z_0)$ real con signo};
  \node[fembox, below=0.8cm of sl2, xshift= 3.5cm] (sl3b)
    {\texttt{slice\_pressure\_field}\\$|p(x,y,z_0,f)|$ complejo $\to|{\cdot}|$};

  \node[fembox, below=0.9cm of sl3a] (sl4a)
    {\texttt{slice\_field\_fem}\\interpolación baricéntrica sobre rejilla $nx\times ny$};
  \node[fembox, below=0.9cm of sl3b] (sl4b)
    {\texttt{modal\_pressure\_field}\\$\to$ \texttt{slice\_field\_fem}};

  \node[vizbox, below=0.9cm of sl4a, xshift=3.5cm] (sl5)
    {\texttt{FieldSliceItem.update()}\\construye quads solo donde\\
     los 4 vértices están \emph{dentro} del recinto};
  \node[vizbox, below=0.8cm of sl5] (sl6)
    {\texttt{GLMeshItem} con \texttt{faceColors}\\montado sobre \texttt{IsoViewer}};

  \draw[arr] (sl1) -- (sl2);
  \draw[arr] (sl2) -| node[label, above left]{\scriptsize kind=0} (sl3a);
  \draw[arr] (sl2) -| node[label, above right]{\scriptsize kind=1} (sl3b);
  \draw[arr] (sl3a) -- (sl4a);
  \draw[arr] (sl3b) -- (sl4b);
  \draw[arr] (sl4a) |- (sl5);
  \draw[arr] (sl4b) |- (sl5);
  \draw[arr] (sl5) -- (sl6);

\end{tikzpicture}
\caption{Flujo del slice 2D horizontal.}
\end{figure}

\paragraph{Enmascaramiento dentro del recinto.}
Un quad de la rejilla se incluye solo si sus cuatro vértices tienen
coordenadas dentro del mesh FEM (el interpolador devuelve un valor finito):

\[
  \text{mostrar quad}_{ij} \;\iff\;
  \mathrm{mask}[i,j] \;\wedge\; \mathrm{mask}[i{+}1,j] \;\wedge\;
  \mathrm{mask}[i,j{+}1] \;\wedge\; \mathrm{mask}[i{+}1,j{+}1]
\]

Esto evita que el slice ``sangre'' fuera del contorno del recinto.

% ------------------------------------------------------------
\section{Flujo 7 — RT60 y amortiguamiento desde materiales}
% ------------------------------------------------------------

\begin{figure}[H]
\centering
\begin{tikzpicture}[node distance=0.85cm and 0cm]

  \node[guibox] (m1) {Usuario elige material\\(combo Piso / Techo / Paredes)};
  \node[matbox, below=of m1] (m2)
        {\texttt{MaterialLibrary} $\leftarrow$ \texttt{materials/*.json}\\
         interpolación log en bandas de octava};
  \node[matbox, below=of m2] (m3)
        {\texttt{classify\_surface\_areas(verts, tris)}\\
         clasifica tris por normal:\\
         piso ($n_z < -0{,}7$) / techo ($n_z > 0{,}7$) / paredes (resto)};
  \node[matbox, below=of m3] (m4)
        {\texttt{compute\_sabine\_rt60(V, Sf, mf, Sc, mc, Sw, mw)}\\
         $A(f) = \alpha_f(f)\,S_f + \alpha_c(f)\,S_c + \alpha_w(f)\,S_w$\\
         $\mathrm{RT60}(f) = 0{,}161\,V/A(f)$};
  \node[matbox, below=of m4] (m5)
        {\texttt{compute\_xi\_per\_mode(freqs, rt60\_bands)}\\
         $\xi_n = 1{,}1 / [f_n \cdot \mathrm{RT60}(f_n)]$};
  \node[fembox, below=of m5] (m6)
        {$\xi_n$ array $(N_m,)$\\pasado a \texttt{frequency\_response} y\\
         \texttt{modal\_pressure\_field}};
  \node[vizbox, below=of m6] (m7)
        {\texttt{RT60PlotDialog}\\gráfico matplotlib de RT60$(f)$\\
         Export PNG/SVG};

  \foreach \i/\j in {m1/m2, m2/m3, m3/m4, m4/m5, m5/m6, m6/m7}
    \draw[arr] (\i) -- (\j);

\end{tikzpicture}
\caption{Flujo de cálculo de amortiguamiento modal desde materiales.}
\end{figure}

% ------------------------------------------------------------
\section{Tabla resumen de señales Qt}
% ------------------------------------------------------------

\begin{table}[H]
\centering
\renewcommand{\arraystretch}{1.3}
\begin{tabular}{>{\ttfamily}l l >{\ttfamily}l}
\toprule
\normalfont\textbf{Señal} & \textbf{Emisor} & \normalfont\textbf{Slot receptor} \\
\midrule
parametersChanged(dict)      & ControlPanel  & MainWindow.\_on\_params \\
parametersCommitted(dict)    & ControlPanel  & MainWindow.\_on\_commit \\
cameraResetRequested()       & ControlPanel  & IsoViewer.reset\_camera \\
drawShapeRequested()         & ControlPanel  & MainWindow.\_open\_shape\_dialog \\
showLabelsToggled(bool)      & ControlPanel  & IsoViewer.set\_show\_labels \\
viewModeChanged(str)         & ControlPanel  & IsoViewer.set\_view\_mode \\
wallDragStarted(int)         & IsoViewer     & ControlPanel.begin\_wall\_drag \\
wallDragMoved(float)         & IsoViewer     & ControlPanel.update\_wall\_drag \\
wallDragEnded()              & IsoViewer     & ControlPanel.end\_wall\_drag \\
sourceAddRequested(x,y,z)    & IsoViewer     & MainWindow.\_on\_source\_add \\
sourceEditRequested(int)     & IsoViewer     & MainWindow.\_on\_source\_edit \\
sourceMoveRequested(i,x,y,z) & IsoViewer     & MainWindow.\_on\_source\_moved \\
receiverMoveRequested(x,y,z) & IsoViewer     & MainWindow.\_on\_receiver\_moved \\
sourceAddedAtFloor(x,y)      & ShapeCanvas   & MainWindow.\_on\_source\_add \\
\bottomrule
\end{tabular}
\caption{Tabla de señales Qt en la aplicación.}
\end{table}

% ------------------------------------------------------------
\section{Ciclo de vida del \texttt{ModalSolution}}
% ------------------------------------------------------------

\begin{figure}[H]
\centering
\begin{tikzpicture}[node distance=1.0cm and 0.5cm,
    state/.style={draw, rounded corners=6pt, minimum width=3.8cm,
                  minimum height=0.8cm, text centered, font=\small}]

  \node[state, fill=clrGUI!40]  (idle)    {Sin cálculo};
  \node[state, fill=clrFEM!50, right=2.5cm of idle]  (calc)
        {Calculando modos};
  \node[state, fill=clrGeo!50, right=2.5cm of calc]  (valid)
        {modal\_result válido};
  \node[state, fill=gray!20,   below=1.5cm of valid] (invalid)
        {modal\_result = None};

  \draw[arr] (idle)    -- node[label, above]{\scriptsize btn\_solve\_fem} (calc);
  \draw[arr] (calc)    -- node[label, above]{\scriptsize eigsh OK} (valid);
  \draw[arr] (calc)    to[bend right=30]
      node[label, below]{\scriptsize error} (idle);
  \draw[arr] (valid)   to[bend left=20]
      node[label, right]{\scriptsize geometría cambia\\
                         (on\_geometry\_changed)} (invalid);
  \draw[arr] (invalid) to[bend left=20]
      node[label, left]{\scriptsize btn\_solve\_fem} (calc);

\end{tikzpicture}
\caption{Ciclo de vida del objeto \texttt{ModalSolution}.
         Cualquier cambio de geometría invalida el resultado.}
\end{figure}

% ------------------------------------------------------------
\section{Consideraciones de implementación}
% ------------------------------------------------------------

\subsection{Orden de renderizado en OpenGL}

Los ítems del \texttt{GLViewWidget} se renderizan en el orden en que fueron
agregados.  Las esferas de fuentes (\texttt{GLScatterPlotItem} con
\texttt{pxMode=True}) deben agregarse \emph{después} del mesh del recinto
para garantizar visibilidad.  Por eso \texttt{\_on\_params()} llama a
\texttt{\_refresh\_sources\_list()} tras \texttt{update\_geometry()}.

\subsection{Proyección manual de la cámara}

pyqtgraph 0.14 cambió la signatura de \texttt{projectionMatrix()}.
La aplicación construye manualmente las matrices view y projection con
\texttt{QMatrix4x4} para el cálculo de etiquetas de aristas y picking 2D,
sin depender de la API interna de pyqtgraph.

\subsection{Debounce del campo acústico}

El arrastre de fuentes emite \texttt{sourceMoveRequested} en cada evento
de mouse (potencialmente 60~veces por segundo).  Un \texttt{QTimer}
monoshot de 350~ms actúa como debounce: el campo se recalcula solo cuando
el usuario detiene el arrastre, evitando saturar la CPU con cómputos FEM
intermedios.

\subsection{Interpolación baricéntrica en mallas tet}

El \texttt{FieldEvaluator} pre-calcula matrices de transformación afín
$A_e^{-1}$ para cada tetraedro.  La localización de un punto $\mathbf{x}$
se realiza buscando el tetraedro con máximo mínimo de coordenadas
baricéntricas $\min(N_0, N_1, N_2, N_3)$, lo que selecciona el elemento
más ``interno'' al punto en caso de múltiples candidatos.

\end{document}
